% titel: Blatt 7, A7: Dichte und Molmasse einer fluorierten Methanverbindung
% skript: Mortimer Kap. 11 (ideale Gase: Dichte) + Kap. 9 (Lewis)

% Teil I: Werkzeuge, kein L\"osungsweg.

Die Dichte eines (idealen) Gases h\"angt mit der Molmasse zusammen. Aus der
Dichte bei gegebenem $p,T$ l\"asst sich $M$ bestimmen und daraus die
Zusammensetzung $\mathrm{CF}_x\mathrm{H}_{4-x}$.

\begin{werkzeugkasten}[Dichte eines idealen Gases]
\[
  \rho = \frac{m}{V} = \frac{pM}{RT}
  \quad\Longrightarrow\quad
  M = \frac{\rho RT}{p}.
\]
Mit $\rho$ in $\mathrm{g\,L^{-1}}$, $p$ in atm, $R=0.0821\ \mathrm{L\,atm\,mol^{-1}K^{-1}}$
folgt $M$ in $\mathrm{g\,mol^{-1}}$. Bei anderem $p,T$ skaliert $\rho
\propto p/T$ (da $M$ konstant).
\textit{Skript:} Mortimer Kap. 11 (Dichte idealer Gase).
\end{werkzeugkasten}

\begin{achtung}[Darauf achten]
\begin{itemize}
  \item Molmasse der Reihe $\mathrm{CF}_x\mathrm{H}_{4-x}$:
        $M = 12.01 + 19.00\,x + 1.008\,(4-x)$.
  \item $x$ m\"ussen ganzzahlig sein; die n\"achste ganze Zahl zur berechneten
        Molmasse liefert $x$.
  \item Bei neuen $p,T$: $\rho_2 = \rho_1\,\frac{p_2}{p_1}\cdot\frac{T_1}{T_2}$.
\end{itemize}
\end{achtung}
